\documentclass[11pt, a4paper]{article}
\usepackage{geometry}
\usepackage[parfill]{parskip}
\usepackage{graphicx}
\usepackage{amsmath}
\usepackage{enumerate}

\title{Assignment 1 \\ IB3702 Mathematics for Machine learning}
\author{Harman Singh} % please overwrite this line with your name
\date{Month Day, Year} % todo please overwrite this line with the date of submission

\begin{document}
\maketitle

\section*{Problem 1}

\textbf{Question 1.1}
\begin{enumerate}[(a)]
	\item Considering $a_0$ is 1, we have to simply fill in 1 where it says $a_n$ in the formula. The result of this will be $a_1$. Then we will take $a_1$ and fill it in the formula to get $a_2$. We will repeat this process until we reach $a_4$.
	\begin{itemize}
		\item $a_{0} = 1$
		\item $a_{1} = \frac{1}{2} (1             	+ \frac{2}{1          })   	= \frac{1}{2} * 3 = \frac{3}{2}$ 																						\hfill \textbf{So $a_1$ = $\frac{3}{2}$}
		\item $a_{2} = \frac{1}{2} (\frac{3}{2}   	+ \frac{2}{\frac{3}{2}})   	= \frac{1}{2} * (\frac{3}{2} + \frac{4}{3}) = \frac{1}{2} * \frac{17}{6} = \frac{17}{12}$ 								\hfill \textbf{So $a_2$ = $\frac{17}{12}$}
		\item $a_{3} = \frac{1}{2} (\frac{17}{12} 	+ \frac{2}{\frac{17}{12}}) 	= \frac{1}{2} * (\frac{17}{12} + \frac{24}{17}) = \frac{1}{2} * \frac{577}{204} = \frac{577}{408} $ 					\hfill \textbf{So $a_3$ = $\frac{577}{408}$}
		\item $a_{4} = \frac{1}{2} (\frac{577}{408} + \frac{2}{\frac{577}{408}})= \frac{1}{2} * (\frac{577}{408} + \frac{816}{577}) = \frac{1}{2} * \frac{665,857}{235,416} = \frac{665,857}{470,832}$ 	\hfill \textbf{So $a_4$ = $\frac{665,857}{470,832}$}
	\end{itemize}
	
	
	
	\item answer/solution to this problem
	

	% We want to prove the limit is √2. The sequence seems positive and likely converges to √2.
	% Using the function 𝑓(𝑥)=1\2(𝑥+2\𝑥), which is related to Newton's method, we need to show
	% the sequence is monotone—either increasing or decreasing based on whether 𝑎_𝑛 is less
	% than or greater than √2.



\end{enumerate}

\section*{Problem 2}

\begin{enumerate}[(a)]
	\item answer/solution to this problem
	\item answer/solution to this problem
\end{enumerate}

\section*{Problem 3}

\begin{enumerate}[(a)]
	\item answer/solution to this problem
	\item answer/solution to this problem
\end{enumerate}

\ldots
\end{document}  